
\section{Round-fourteen positive closure: a positive phase sheaf and logarithmic nonlinear mixtures}
\label{sec:r14-d1}

The phase label is a positive measurable coordinate on the original
microscopic sample space.  Component rates are lower-semicontinuous
basin-internal recovery costs, not the unrestricted rate on a basin boundary.
Nonlinear log-Laplace values combine by an exact log-sum-exp formula, never
by a linear mixture.

\subsection{Canonical rate basins}

Let \(X_n\) be a platform path variable with good rate \(I\), and let
\(M:\mathcal X\to(\mathcal O,d_{\mathcal O})\) be the platform order parameter
fixed in A3 or B2.  On a regular finite-phase chart assume the zero/low-rate
set of \(I_M\) has finitely many isolated connected components
\(K_1,\ldots,K_J\).  Choose \(\delta\) below one quarter of their mutual
distance and define the intrinsic Voronoi basins
\[
 G_j=\{x:d(Mx,K_j)<\delta,\ 
          d(Mx,K_j)<d(Mx,K_\ell),\ \ell\ne j\},
\]
and \(G_\partial=\mathcal X\setminus\bigcup_jG_j\).
Ties are assigned to the boundary label.  This construction depends only on
the metric rate components and is invariant under isometric
reparametrization of the order parameter.

Let \(J_n=J(X_n)\in\{1,\ldots,J,\partial\}\).

\begin{theorem}[Exact positive phase sheaf]
\label{thm:r14-d1-label}
For every \(n\),
\[
 \mathbb P_n
 =
 \sum_{j\in\{1,\ldots,J,\partial\}}
 w_{n,j}\mathbb P_{n,j},
 \qquad
 w_{n,j}=\mathbb P_n(J_n=j),
\]
where every \(\mathbb P_{n,j}\) is a probability conditional on the same
measurable label.  The label is unchanged under every finite observable
projection, so it defines one projective labelled process.  The boundary
component is retained at its actual cost.
\end{theorem}

\begin{proof}
The intrinsic Voronoi rule is Borel and gives a disjoint partition.  The
identity is the law of total probability.  Since the label is defined once
from the full order parameter, all projected laws carry the same random
label; no projection-dependent relabelling occurs.
\end{proof}

\subsection{Component rates with boundary recovery}

Define the basin-interior regularization
\[
 I_j(x)=
 \lim_{\epsilon\downarrow0}
 \inf\{I(y)-\alpha_j:
       y\in G_j,\ d(x,y)<\epsilon\},
 \qquad
 \alpha_j=\inf_{y\in G_j}I(y),
\]
with value \(+\infty\) if the set is empty.  Define \(I_\partial\) in the same
way.

\begin{lemma}[Basin-internal recovery]
\label{lem:r14-d1-recovery}
For the A3 and B2 platform rates, every \(x\) with \(I_j(x)<\infty\) has a
microscopic recovery sequence whose empirical state remains in \(G_j\) and
whose cost tends to \(\alpha_j+I_j(x)\).
\end{lemma}

\begin{proof}
By definition choose \(y_m\in G_j\) tending to \(x\) with
\(I(y_m)\to\alpha_j+I_j(x)\).  A3 recovers each finite-rate controlled renewal
law by an actual stationary controlled kernel and legal connectors.  B2
recovers each finite-action density/contact path by a smooth positive exposed
pair and the exact kinetic Hodge repair.  Because \(G_j\) is open and
\(M\) is continuous, the recovery may be chosen in a sufficiently small
neighborhood of \(y_m\), hence inside \(G_j\).  Diagonalizing in \(m\) gives
the sequence.  The same argument applies to the boundary component using its
own interior in the labelled topology.
\end{proof}

\begin{theorem}[Positive component LDP]
\label{thm:r14-d1-component}
If \(-n^{-1}\log w_{n,j}\to\alpha_j\), the conditional laws
\(\mathbb P_{n,j}\) satisfy good LDPs with rates \(I_j\).  The joint labelled
law has rate
\[
 \mathcal J(j,x)=\alpha_j+I_j(x),
\]
and forgetting the label gives
\[
 I(x)=\min_j\{\alpha_j+I_j(x)\}.
\]
At a basin boundary, each component pays its own basin-internal recovery
cost; it is not automatically \(I(x)-\alpha_j\).
\end{theorem}

\begin{proof}
For a closed set use the platform upper bound on its intersection with the
closed basin closure.  For an open set and a finite-rate point, use
Lemma~\ref{lem:r14-d1-recovery}.  Division by \(w_{n,j}\) subtracts
\(\alpha_j\).  Compactness follows because component sublevels are closed
subsets of a compact \(I\)-sublevel.  The finite-label joint LDP and the
contraction principle give the final formula.
\end{proof}

\subsection{Component source charts}

Let
\[
 \widetilde Q_{n,j}(F)
 ={1\over n}\log
 \mathbb E_n[e^{nF(X_n)}\mid J_n=j].
\]
The unnormalized restricted pressure is
\[
 Q_{n,j}^{\rm un}(F)
 ={1\over n}\log
 \mathbb E_n[e^{nF(X_n)};J_n=j]
 =-{1\over n}\log w_{n,j}+\widetilde Q_{n,j}(F).
\]

\begin{theorem}[Phase-conditioned analytic chart]
\label{thm:r14-d1-chart}
If the tilted minimizer of \(\alpha_j+I_j-\nu(F)\) is unique and lies a
positive rate distance from the labelled boundary, then
\(\widetilde Q_{n,j}\) has a zero-free complex neighborhood with locally
uniform first two derivatives.  Its limiting derivatives are the component
mean and covariance.  The boundary label has the same theorem whenever it
has an interior exposed phase; it is otherwise retained through its real
Laplace principle without being declared negligible.
\end{theorem}

\begin{proof}
For A3, impose the basin label on the controlled Markov-renewal kernel.
The positive killed transfer operator has leakage to the other labelled
components of size \(e^{-gn}\), where \(g\) is the rate separation.  Its
Perron root is isolated on the component graph, and analytic perturbation in
the source gives a zero-free coefficient chart.

For B2, the component condition is an initial/order-parameter constraint.
B1 exact-number coefficient extraction and the source-dependent saddle apply
with the positive basin indicator inserted.  The interior rate gap makes the
complement exponentially smaller relative to the component coefficient.
In both cases Rouch\'e's theorem gives the finite-volume zero-free
neighborhood and Cauchy estimates give derivative convergence.  No real
probability gap alone is used without the corresponding positive component
operator/coefficient construction.
\end{proof}

\subsection{Exact nonlinear phase aggregation}

Let \(V_{n,j}(\Phi)\) be any component log-Laplace value, including a dynamic
component Nisio/Doob value.  Then exactly
\[
 V_n(\Phi)
 ={1\over n}\log
 \sum_j
 \exp\{n[-\alpha_{n,j}+V_{n,j}(\Phi)]\},
 \qquad
 \alpha_{n,j}=-{1\over n}\log w_{n,j}.
\]

\begin{theorem}[Logarithmic phase-mixture calculus]
\label{thm:r14-d1-logsum}
If the component values converge locally uniformly, then
\[
 V(\Phi)=\max_j\{-\alpha_j+V_j(\Phi)\}.
\]
At a unique maximizer, derivatives equal that component's derivatives.  At a
tie, the subdifferential is the convex hull of the active component
derivatives, with finite-volume subleading weights selecting subsequential
mixtures.  Linear evolution acts on component measures before the logarithm;
the nonlinear values themselves are never averaged linearly.
\end{theorem}

\begin{proof}
The finite formula is the exact positive disintegration from
Theorem~\ref{thm:r14-d1-label}.  The elementary Laplace principle for a finite
sum gives the maximum.  Differentiate the finite log-sum-exp: the derivative
is the softmax-weighted component derivative.  At a unique maximizer the
weight converges to one; at a tie every cluster point lies in the stated
convex hull.  This proves the nonlinear aggregation rule.
\end{proof}

For dynamic values, retain the finite phase weights as a state coordinate.
Component semigroups evolve each positive component, and Bayes/log-sum-exp
updates the weights.  This gives an augmented nonlinear semigroup even at
coexistence.

\subsection{Phase-conditioned shells and tangents}

Let \(C(X_n)\) be a conditioned macro variable.  Suppose each active component
has target rate
\[
 I_{j,C}(a)=\inf_{Cx=a}I_j(x)
\]
and local coefficient \(c_{n,j}(a)\).

\begin{theorem}[Exact target-dependent shell mixture]
\label{thm:r14-d1-shell}
The phase weight after conditioning on a regular shell about \(a\) is
proportional to
\[
 w_{n,j}\,c_{n,j}(a)\,e^{-nI_{j,C}(a)}.
\]
The boundary label is included with its own target rate and coefficient.
A uniquely smallest exponent selects one phase; at a tie, the normalized
subexponential coefficients determine subsequential mixtures.  The Gaussian
tangent in a selected phase is its component covariance, while at a tie the
labelled tangent is a finite Gaussian mixture rather than one artificially
averaged Gaussian.
\end{theorem}

\begin{proof}
Intersect the exact positive disintegration with the shell.  Apply the A2 or
B1 component coefficient theorem to each positive term.  Each term has its
own constrained infimum \(I_{j,C}(a)\), including basin-boundary recovery.
Divide by their sum.  Conditional Gaussian limits pass termwise because the
number of labels is finite.
\end{proof}

\subsection{Commutation on the labelled law}

\begin{theorem}[Positive labelled commutation theorem]
\label{thm:r14-d1-commute}
On the joint labelled microscopic law, measurable contraction, exact shell
conditioning, component Gaussian scaling, and linear law evolution commute
with positive phase disintegration.  After applying the logarithm, their
nonlinear aggregation is precisely
Theorem~\ref{thm:r14-d1-logsum}; no linear-mixture commutation is asserted.
\end{theorem}

\begin{proof}
Before the logarithm, all operations act on the finite positive measures
\(w_{n,j}\mathbb P_{n,j}\): pushforward, restriction, rescaling, and linear
law evolution distribute over a finite sum.  Conditioning renormalizes the
resulting positive terms.  Taking the logarithm after summation gives the
exact log-sum-exp identity and no other aggregation rule.
\end{proof}

\begin{theorem}[Round-fourteen D1 closure]
\label{thm:r14-d1-main}
Every regular finite-phase chart of the proved Sinai or hard-sphere platform
has a projectively consistent positive phase sheaf with basin-internal
component rates, component response charts, exact target-dependent shell
weights, Gaussian mixtures, and logarithmic nonlinear dynamic aggregation.
The phase cost is counted once and no spectral projection is interpreted as
a probability.
\end{theorem}

\begin{proof}
Combine Theorems~\ref{thm:r14-d1-label},
\ref{thm:r14-d1-component}, \ref{thm:r14-d1-chart},
\ref{thm:r14-d1-logsum}, \ref{thm:r14-d1-shell}, and
\ref{thm:r14-d1-commute}.
\end{proof}
